% !TeX program = pdflatex

%\documentclass[ngerman, english]{udeddi} % 4:3 aspec ratio
\documentclass[aspectratio=169, ngerman]{udeddi}

\usepackage[babel,german=quotes]{csquotes}
\usepackage[style=authoryear-comp,natbib=true,backend=biber,maxcitenames=2,maxbibnames=2]{biblatex}

\usepackage[mode=buildnew]{standalone}
\usepackage{booktabs}

\bibliography{presentation.bib}
\AtEveryBibitem{%
    \clearfield{issn} % Remove issn
    \clearfield{doi} % Remove doi
    \clearfield{eprint}
    
    \ifentrytype{online}{}{% Remove url except for @online
        \clearfield{url}
        \clearfield{urldate}
        \clearfield{urlmonth}
        \clearfield{urlyear}
    }
}

% use to show notes
% \usepackage{pgfpages}
% \setbeameroption{show notes on second screen=right}

\title{Titel der Präsentation}
\topic{Untertitel der Präsentation oder Veranstaltung}
\footertitle{Kurztitel}

\author{Mike Barkmin}
% \author{Mike Barkmin, Torsten Brinda \& Matthias Kramer}
\date{07. Oktober 2020}
\mail{mike.barkmin@uni-due.de}
% \mail{\{mike.barkmin, torsten.brinda, matthias.kramer\}@uni-due.de}
\website{http://udue.de/mba} % optional
% \website{http://udue.de/\{mba, tba, mk\}}
\organisation{Universität Duisburg-Essen}
\chair{Didaktik der Informatik}
\address{Schützenbahn 70, 45127 Essen}
\slidesurl{https://www.ddi.wiwi.uni-due.de/}

\begin{document}
\maketitle

\begin{frame}{Inhaltsverzeichnis}
  \maintableofcontents
\end{frame}

\section{Zitate}
\begin{frame}{Definition Kompetenz \citep{weinert_vergleichende_2001}}
    \begin{quote}
        \enquote{Dabei versteht man unter Kompetenzen die bei Individuen verfügbaren
            oder durch sie erlernbaren kognitiven Fähigkeiten und Fertigkeiten,
            um bestimmte Probleme zu lösen, sowie die damit verbundenen
            motivationalen, volitionalen und sozialen Bereitschaften und
            Fähigkeiten um die Problemlösungen in variablen Situationen
            erflogreich und verantwortungsvoll nutzen zu können.}
    \end{quote}
\end{frame}

\section{Zwei Spalten}
\begin{frame}{Zielstellungen}
    \begin{columns}
        \begin{column}{0.5\textwidth}
            \begin{itemize}
                \item Ziel 1.1
                \item Ziel 1.2
                \item Ziel 1.3
                \item Ziel 1.4
            \end{itemize}
        \end{column}
        \begin{column}{0.5\textwidth}
            \begin{itemize}
                \item Ziel 2.1
                \item Ziel 2.2
                \item Ziel 2.3
                \item Ziel 2.4
            \end{itemize}
        \end{column}
    \end{columns}
\end{frame}

\section{Smarte Diagramme}
\begin{frame}{Circular - Primary}
\begin{center}
    \smartdiagram[circular diagram]{Set up,Run,Analyse,Modify~/ Add}
\end{center}
\end{frame}

\begin{frame}{Circular - Secondary}
\begin{center}
    \smartdiagramset{secondary}
    \smartdiagram[circular diagram]{Set up,Run,Analyse,Modify~/ Add}
\end{center}
\end{frame}

\begin{frame}{Circular - Color}
\begin{center}
    \smartdiagramset{color}
    \smartdiagram[circular diagram]{Set up,Run,Analyse,Modify~/ Add}
\end{center}
\end{frame}

\begin{frame}{Sequence - Primary}
\begin{center}
    \smartdiagramset{primary}
    \smartdiagram[sequence diagram]{PGF,Ti\textit{k}Z,Smartdiagram}
\end{center}
\end{frame}

\begin{frame}{Sequence - Secondary}
\begin{center}
    \smartdiagramset{secondary}
    \smartdiagram[sequence diagram]{PGF,Ti\textit{k}Z,Smartdiagram}
\end{center}
\end{frame}

\begin{frame}{Sequence - Color}
\begin{center}
    \smartdiagramset{color}
    \smartdiagram[sequence diagram]{PGF,Ti\textit{k}Z,Smartdiagram}
\end{center}
\end{frame}

\begin{frame}{More Diagrams}
    \begin{center}
        Visit the documentation of smartdiagram for more: 
        \url{http://texdoc.net/texmf-dist/doc/latex/smartdiagram/smartdiagram.pdf}
    \end{center}
\end{frame}

\begin{contactFrame}{\insertauthor}{\insertchair}{\insertorganisation}{\insertaddress}{\insertmail}{\insertwebsite}
    Vielen Dank für die Aufmerksamkeit!
\end{contactFrame}

\appendix

\begin{bibFrame}[Literaturverzeichnis]
    \printbibliography[heading=none]
\end{bibFrame}

\end{document}
